% problem_id: S001-lemma-from-functor-to-weil
% source_id: S001 (Stacks Project)
% source_file: weil.tex
% statement_locator: weil.tex lines 2792--2802
% proof_locator: weil.tex lines 2804--2977
% env: lemma
% id_note: label 在本来源内唯一
% pairing: tight (verified)
% status: complete (statement + author proof verbatim + reason + locators)
%
\begin{lemma}
\label{lemma-from-functor-to-weil}
Let $k$ be a field. Let $F$ be a field of characteristic $0$.
Assume given a $\mathbf{Q}$-linear functor
$$
G : M_k \longrightarrow \text{graded }F\text{-vector spaces}
$$
of symmetric monoidal categories such that $G(\mathbf{1}(1))$
is nonzero only in degree $-2$. Then we obtain data (D0), (D1), (D2), and (D3)
satisfying all of (A), (B), and (C) above.
\end{lemma}

\begin{proof}
This proof is the same as the proof of
Lemma \ref{lemma-from-functor-to-weil-classical};
we urge the reader to read the proof of that lemma instead.

\medskip\noindent
We obtain a contravariant functor from the category of smooth
projective schemes over $k$ to the category of graded $F$-vector spaces
by setting $H^*(X) = G(h(X))$. By assumption we have a canonical
isomorphism
$$
H^*(X \times Y) = G(h(X \times Y)) = G(h(X) \otimes h(Y)) =
G(h(X)) \otimes G(h(Y)) = H^*(X) \otimes H^*(Y)
$$
compatible with pullbacks. Using pullback along the diagonal
$\Delta : X \to X \times X$ we obtain a canonical map
$$
H^*(X) \otimes H^*(X) = H^*(X \times X) \to H^*(X)
$$
of graded vector spaces compatible with pullbacks.
This defines a functorial graded $F$-algebra structure on
$H^*(X)$. Since $\Delta$ commutes with the commutativity
constraint $h(X) \otimes h(X) \to h(X) \otimes h(X)$ (switching the factors)
and since $G$ is a functor of symmetric monoidal categories (so compatible with
commutativity constraints), and by our convention in
Homology, Example \ref{homology-example-graded-vector-spaces}
we conclude that $H^*(X)$ is a graded
commutative algebra! Hence we get our datum (D1).

\medskip\noindent
Since $\mathbf{1}(1)$ is invertible in the category of motives
we see that $G(\mathbf{1}(1))$ is invertible in the category of
graded $F$-vector spaces. Thus $\sum_i \dim_F G^i(\mathbf{1}(1)) = 1$.
By assumption we only get something nonzero in degree $-2$.
Our datum (D0) is the vector space $F(1) = G^{-2}(\mathbf{1}(1))$.
Since $G$ is a symmetric monoidal functor we see that
$F(n) = G^{-2n}(\mathbf{1}(n))$ for all $n \in \mathbf{Z}$.
It follows that
$$
H^{2r}(X)(r) = G^{2r}(h(X)) \otimes G^{-2r}(\mathbf{1}(r)) =
G^0(h(X)(r))
$$
a formula we will frequently use below.

\medskip\noindent
Let $X$ be a smooth projective scheme over $k$. By
Lemma \ref{lemma-composition-correspondences} we have
$$
\CH^r(X) \otimes \mathbf{Q} = \text{Corr}^r(\Spec(k), X) =
\Hom(\mathbf{1}(-r), h(X)) = \Hom(\mathbf{1}, h(X)(r))
$$
Applying the functor $G$ this maps into
$\Hom(G(\mathbf{1}), G(h(X)(r)))$.
By taking the image of $1$ in $G^0(\mathbf{1}) = F$
into $G^0(h(X)(r)) = H^{2r}(X)(r)$ we obtain
$$
\gamma :
\CH^r(X) \otimes \mathbf{Q} \longrightarrow H^{2r}(X)(r)
$$
This is the datum (D2).

\medskip\noindent
Let $X$ be a nonempty smooth projective scheme over $k$
which is equidimensional of dimension $d$. By
Lemma \ref{lemma-composition-correspondences} we have
$$
\Mor(h(X)(d), \mathbf{1}) = \Mor((X, 1, d), (\Spec(k), 1, 0)) =
\text{Corr}^{-d}(X, \Spec(k)) = \CH_d(X)
$$
Thus the class of the cycle $[X]$ in $\CH_d(X)$ defines a morphism
$h(X)(d) \to \mathbf{1}$. Applying $G$ and taking degree $0$
parts we obtain
$$
H^{2d}(X)(d) = G^0(h(X)(d)) \longrightarrow G^0(\mathbf{1}) = F
$$
This map $\int_X : H^{2d}(X)(d) \to F$ is the datum (D3).

\medskip\noindent
Let $X$ be a smooth projective scheme over $k$ which is
nonempty and equidimensional of dimension $d$. By Lemma \ref{lemma-dual}
we know that $h(X)(d)$ is a left dual to $h(X)$. Hence
$G(h(X)(d)) = H^*(X) \otimes_F F(d)[2d]$
is a left dual to $H^*(X)$ in the category of graded $F$-vector spaces.
Here $[n]$ is the shift functor on graded vector spaces.
By Homology, Lemma \ref{homology-lemma-left-dual-graded-vector-spaces}
we find that $\sum_i \dim_F H^i(X) < \infty$ and that
$\epsilon : h(X)(d) \otimes h(X) \to \mathbf{1}$ produces
nondegenerate pairings $H^{2d - i}(X)(d) \otimes_F H^i(X) \to F$.
In the proof of Lemma \ref{lemma-dual} we have seen that
$\epsilon$ is given by $[\Delta]$ via the identifications
$$
\Hom(h(X)(d) \otimes h(X), \mathbf{1}) =
\text{Corr}^{-d}(X \times X, \Spec(k)) =
\CH_d(X \times X)
$$
Thus $\epsilon$ is the composition of $[X] : h(X)(d) \to \mathbf{1}$
and $h(\Delta)(d) : h(X)(d) \otimes h(X) \to h(X)(d)$. It follows
that the pairings above are given by cup product followed by
$\int_X$. This proves axiom (A).

\medskip\noindent
Axiom (B) follows from the assumption that $G$ is compatible
with tensor structures and our construction of the cup product above.

\medskip\noindent
Axiom (C). Our construction of $\gamma$ takes a cycle $\alpha$ on $X$,
interprets it a correspondence $a$ from $\Spec(k)$ to $X$ of some degree,
and then applies $G$. If $f : Y \to X$ is a morphism of nonempty
equidimensional smooth projective schemes over $k$, then
$f^!\alpha$ is the pushforward (!) of $\alpha$
by the correspondence $[\Gamma_f]$ from $X$ to $Y$, see
Lemma \ref{lemma-functor-and-cycles}. Hence
$f^!\alpha$ viewed as a correspondence from $\Spec(k)$ to $Y$
is equal to $a \circ [\Gamma_f]$, see
Lemma \ref{lemma-composition-correspondences}.
Since $G$ is a functor, we conclude
$\gamma$ is compatible with pullbacks, i.e., axiom (C)(a) holds.

\medskip\noindent
Let $f : Y \to X$ be a morphism of nonempty equidimensional
smooth projective schemes over $k$ and
let $\beta \in \CH^r(Y)$ be a cycle on $Y$. We have to show that
$$
\int_Y \gamma(\beta) \cup f^*c = \int_X \gamma(f_*\beta) \cup c
$$
for all $c \in H^*(X)$. Let $a, a^t, \eta_X, \eta_Y, [X], [Y]$
be as in Lemma \ref{lemma-prep-dual}.
Let $b$ be $\beta$ viewed as a correspondence from $\Spec(k)$ to $Y$
of degree $r$. Then $f_*\beta$ viewed as a correspondence from
$\Spec(k)$ to $X$ is equal to $a^t \circ b$, see
Lemmas \ref{lemma-functor-and-cycles} and
\ref{lemma-composition-correspondences}.
The displayed equality above holds if we can show that
$$
h(X) = \mathbf{1} \otimes h(X)
\xrightarrow{b \otimes 1}
h(Y)(r) \otimes h(X)
\xrightarrow{1 \otimes a}
h(Y)(r) \otimes h(Y)
\xrightarrow{\eta_Y}
h(Y)(r)
\xrightarrow{[Y]}
\mathbf{1}(r - e)
$$
is equal to
$$
h(X) = \mathbf{1} \otimes h(X)
\xrightarrow{a^t \circ b \otimes 1}
h(X)(r + d - e) \otimes h(X)
\xrightarrow{\eta_X}
h(X)(r + d - e)
\xrightarrow{[X]}
\mathbf{1}(r - e)
$$
This follows immediately from Lemma \ref{lemma-prep-dual}.
Thus we have axiom (C)(b).

\medskip\noindent
To prove axiom (C)(c) we use the discussion in
Remark \ref{remark-replace-cup-product-classical}.
Hence it suffices to prove that $\gamma$ is compatible with
exterior products. Let $X$, $Y$ be nonempty smooth projective
schemes over $k$ and let $\alpha$, $\beta$ be cycles on them. Denote
$a$, $b$ the corresponding correspondences from $\Spec(k)$ to
$X$, $Y$. Then $\alpha \times \beta$ corresponds to the
correspondence $a \otimes b$ from $\Spec(k)$ to $X \otimes Y = X \times Y$.
Hence the requirement follows from the fact that $G$ is
compatible with the tensor structures on both sides.

\medskip\noindent
Axiom (C)(d) follows because the cycle $[\Spec(k)]$
corresponds to the identity morphism on $h(\Spec(k))$.
This finishes the proof of the lemma.
\end{proof}
